(1) 在惯量主轴系中无矩刚体满足 $I_1\dot\omega_1=(I_2-I_3)\omega_2\omega_3$、$I_2\dot\omega_2=(I_3-I_1)\omega_3\omega_1$、$I_3\dot\omega_3=(I_1-I_2)\omega_1\omega_2$。代入 $I_1=I_2=I$，得 $I\dot\omega_1=(I-J)\omega_2\omega_3$、$I\dot\omega_2=(J-I)\omega_1\omega_3$、$J\dot\omega_3=0$。
(2) 由第三式 $\omega_3=\Omega_3$ 为常量。置 $w=\omega_1+i\omega_2$，前两式合为 $\dot w=i[(J-I)\Omega_3/I]w$，因此 $w(t)=w(0)\exp(i\nu t)$，$\nu=(J-I)\Omega_3/I$。横向角速度在体轴系绕对称轴均匀旋转，$|w|$ 不变，$\nu$ 的符号由长扁或扁长的转动惯量排列决定。
(3) 动能 $K=\tfrac12I(\omega_1^2+\omega_2^2)+\tfrac12J\Omega_3^2$，由 $|w|$ 恒定可知其为常数。角动量体坐标分量为 $(I\omega_1,I\omega_2,J\Omega_3)$，故 $|\boldsymbol L|^2=I^2|w|^2+J^2\Omega_3^2$ 亦恒定。尽管其体坐标分量旋转，空间惯性系中的角动量矢量不随时间改变；二者不矛盾。
(4) 若 $J=I$，则 $\nu=0$，球对称刚体的体坐标角速度固定；若 $\Omega_3=0$ 亦有 $\nu=0$，横向角速度不旋转。空间中对称轴与恒定角动量的夹角固定，通常绕 $\boldsymbol L$ 以 $|\boldsymbol L|/I$ 的角速率进动（非退化、惯性轴有定义时）；此空间进动率一般不同于体坐标中的 $\nu$，不能把它们等同。
